\documentclass[11pt,a4paper]{article}

\usepackage[spanish,es-tabla,es-nodecimaldot]{babel}
\usepackage[T1]{fontenc}
\usepackage[utf8]{inputenc}
\usepackage{newtxtext,newtxmath}
\usepackage{microtype}
\usepackage[top=2.6cm,bottom=2.5cm,left=2.4cm,right=2.4cm,headheight=26pt]{geometry}
\usepackage{setspace}
\usepackage{graphicx}
\usepackage{xcolor}
\usepackage{booktabs}
\usepackage{tabularx}
\usepackage{array}
\usepackage{float}
\usepackage{enumitem}
\usepackage{fancyhdr}
\usepackage{titlesec}
\usepackage{tocloft}
\usepackage{caption}
\usepackage{colortbl}
\usepackage{tikz}
\usetikzlibrary{shapes.geometric,arrows.meta}
\usepackage{pdflscape}
\usepackage{xurl}
\usepackage[hidelinks]{hyperref}

\definecolor{bcp}{HTML}{0B2545}
\definecolor{cobre}{HTML}{C65D12}
\definecolor{tinta}{HTML}{1C2833}
\definecolor{suave}{HTML}{F3EFE7}

\setstretch{1.15}
\setlength{\parskip}{0.55em}
\setlength{\parindent}{0pt}
\renewcommand{\arraystretch}{1.28}
\captionsetup{font=small,labelfont={bf,color=bcp},skip=6pt,justification=centering}
\captionsetup[table]{position=top}

\pagestyle{fancy}
\fancyhf{}
\fancyhead[L]{\textcolor{bcp}{\small Agente financiero personal}}
\fancyhead[R]{\textcolor{bcp}{\small Línea base de alcance, cronograma y costo}}
\fancyfoot[L]{\textcolor{tinta}{\footnotesize Banca Personas · BCP}}
\fancyfoot[R]{\textcolor{tinta}{\footnotesize\thepage}}
\renewcommand{\headrulewidth}{0.4pt}
\renewcommand{\headrule}{\hbox to\headwidth{\color{cobre}\leaders\hrule height \headrulewidth\hfill}}
\renewcommand{\footrulewidth}{0pt}

\titleformat{\section}{\Large\bfseries\color{bcp}}{\thesection.}{0.6em}{}[\vspace{-0.2em}{\color{cobre}\titlerule[0.8pt]}]
\titleformat{\subsection}{\large\bfseries\color{bcp}}{\thesubsection}{0.6em}{}
\titlespacing*{\section}{0pt}{0.7em}{0.7em}
\titlespacing*{\subsection}{0pt}{0.9em}{0.4em}
\renewcommand{\cftsecleader}{\cftdotfill{\cftdotsep}}
\renewcommand{\cfttableader}{\cftdotfill{\cftdotsep}}
\renewcommand{\cftfigleader}{\cftdotfill{\cftdotsep}}
\setlength{\cftbeforesecskip}{0.5em}
\setlength{\cftbeforetabskip}{0.35em}
\setlength{\cftbeforefigskip}{0.35em}
\setlist[itemize]{leftmargin=1.3em,itemsep=0.25em,topsep=0.3em}
\setlist[enumerate]{leftmargin=1.5em,itemsep=0.25em,topsep=0.3em}

\hypersetup{pdftitle={Línea base de alcance, cronograma y costo — Agente financiero personal},pdfauthor={Franco Stywar Antony Vargas Chino; Albert Christian Hidalgo Sánchez; Sergio Perez Ruiz; Carlos E. Ruiz Perez; Dennis Jack Paitan Cano}}
\newcommand{\usd}[1]{USD~#1}
\newcommand{\clave}[1]{\par\smallskip\noindent\colorbox{suave}{\parbox{\dimexpr\textwidth-2\fboxsep\relax}{\setstretch{1.1}\small #1}}\par\smallskip}

\begin{document}
\pagenumbering{roman}

\begin{center}
\vspace*{0.15cm}
\includegraphics[width=2.7cm]{../images/logo-unu.png}\\[0.55em]
{\color{bcp}\large\bfseries UNIVERSIDAD NACIONAL DE UCAYALI}\\[0.28em]
{\color{bcp} Escuela de Posgrado}\\[0.12em]
Maestría en Ingeniería de Sistemas\\[0.08em]
con mención en Gestión de Tecnologías de la Información\\[0.75em]
Docente: Dr.\ Horacio Soriano Alava, PMP\\[1.05cm]
{\color{cobre}\small Línea base del proyecto}\\[0.55em]
{\color{bcp}\fontsize{20}{25}\selectfont\bfseries Línea base de alcance,\\[0.12em] cronograma y costo}\\[0.5em]
{\Large Agente financiero personal}\\[0.15em]
{\large Banca Móvil y banca por internet}\\[0.35em]
{\color{tinta} Banco de Crédito del Perú · Banca Personas}\\[0.75cm]
{\color{bcp}\rule{0.42\textwidth}{0.7pt}}\\[0.7em]
\begin{tabular}{@{}>{\bfseries\color{bcp}}l l@{}}
Versión & 1.0 \\
Fecha & Noviembre 2026 \\
Estado & Línea base \\
\end{tabular}\\[0.95cm]
{\color{bcp}\bfseries Elaborado por}\\[0.35em]
Franco Stywar Antony Vargas Chino\\
Albert Christian Hidalgo Sánchez\\
Sergio Perez Ruiz\\
Carlos E. Ruiz Perez\\
Dennis Jack Paitan Cano
\vfill
{\small Uso interno}
\end{center}
\thispagestyle{empty}
\clearpage

{\setstretch{1.08}
\tableofcontents
\vspace{1.1em}
\listoftables
\vspace{1.1em}
\listoffigures}
\clearpage
\pagenumbering{arabic}

\section{Objeto}

El acta de constitución autoriza un agente que informa y simula, en Banca Móvil y en la banca por internet. La inversión es \usd{1\,633\,000} y el cierre es el mes~11. Este documento fija la línea base de esa autorización: requisitos y EDT, cronograma con hitos y ruta crítica, y el costo de los mismos paquetes.

Cada paquete de la EDT tiene una semana de aceptación y un monto. La suma es el presupuesto autorizado. No se incorpora ningún monto nuevo.

El libro de trabajo que acompaña este documento calcula fechas, holgura y gasto mensual. La duración y el costo de cada actividad son datos de entrada. Fechas, holgura y totales son fórmulas. El control cierra en cero cuando el CAPEX base es \usd{1\,420\,000} y el presupuesto es \usd{1\,633\,000}.

\section{Requisitos priorizados}

La prioridad sigue MoSCoW. Must es condición para dar por cumplido el proyecto. Should está en la línea base y puede retrasarse sin mover la ruta crítica. Won't queda fuera de la EDT: el acta ya lo excluyó.

\begin{table}[H]
\centering
\caption{Requisitos de la línea base. El costo es el del paquete, compartido por los requisitos que caen en él.}
\footnotesize
\setlength{\tabcolsep}{4pt}
\begin{tabularx}{\textwidth}{@{}l l X l l@{}}
\toprule
\textbf{Id} & \textbf{Prioridad} & \textbf{Requisito} & \textbf{EDT} & \textbf{Hito} \\
\midrule
RQ01 & Must & Toda cifra sale de una vista vigente. Si no lo está, se dice que no está confirmada. & 1.1 & H2 \\
RQ02 & Must & Consentimiento expreso, diferenciado y revocable. Revocar corta la llamada siguiente. & 1.2 & H3 \\
RQ03 & Must & La simulación usa el simulador del producto: cuota, plazo, TCEA y total. & 1.2 & H4 \\
RQ04 & Must & El agente no ejecuta. Confirma la persona y ejecuta el producto. & 1.2 & H5 \\
RQ05 & Must & Hay una vía humana en todo momento. & 1.5 & H5 \\
RQ06 & Must & Cada interacción queda trazada. El modelo no se entrena con la conversación. & 1.5 & H5 \\
RQ07 & Must & Piloto del perfil Consulta con 5\,000 clientes. & 1.3 & H3 \\
RQ08 & Must & Planificación queda en la aplicación y en la web. & 1.3 & H4 \\
RQ09 & Must & Ofertas con confirmación queda en producción, con vía humana. & 1.3 & H5 \\
RQ10 & Must & Los doce casos de comportamiento pasan, sin efecto sobre saldos, pagos ni desembolsos, antes del piloto y antes de producción. & 1.5 & H5 \\
\rowcolor{suave} RQ11 & Should & Descubrimiento con Yape y decisión registrada. No se construye. & 1.4 & H6 \\
\rowcolor{suave} RQ12 & Should & El personal de canales queda habilitado para el piloto y para la salida. & 1.6 & H5 \\
RQ13 & Won't & No asesora, no recomienda valores y no reemplaza a Clara ni a Salvador. & --- & --- \\
RQ14 & Won't & No se rediseña la app, la web ni Yape, y no se escribe en los sistemas de producto. & --- & --- \\
\bottomrule
\end{tabularx}
\end{table}

RQ11 es el único requisito de la línea base que no está en la ruta crítica: el descubrimiento puede terminar hasta la semana~40 sin mover el cierre. RQ12 habilita a las personas; la actividad del piloto tiene cuatro semanas de holgura, y la de la salida a producción está en la ruta porque comparte la ventana con las herramientas de Ofertas.

Cada paquete cubre los requisitos que se listan a continuación.

\begin{table}[H]
\centering
\caption{Matriz de trazabilidad: paquete EDT y requisitos que cubre.}
\small
\begin{tabularx}{\textwidth}{@{}l X l@{}}
\toprule
\textbf{EDT} & \textbf{Requisitos cubiertos} & \textbf{MoSCoW} \\
\midrule
1.1 & RQ01 — Vista vigente; cifra no confirmada si no lo está. & Must \\
1.2 & RQ02 Consentimiento, RQ03 Simulación, RQ04 No ejecutar. & Must \\
1.3 & RQ07 Piloto Consulta, RQ08 Planificación, RQ09 Ofertas. & Must \\
1.4 & RQ11 Descubrimiento Yape. & Should \\
1.5 & RQ05 Vía humana, RQ06 Trazas, RQ10 Doce casos. & Must \\
1.6 & RQ12 Personal habilitado. & Should \\
1.0 & Dirección del proyecto. No apunta a un requisito funcional. & — \\
\bottomrule
\end{tabularx}
\end{table}

RQ13 y RQ14 no tienen paquete. Están en la lista para dejar constancia de que no entran al alcance.

\section{EDT y diccionario}

La EDT tiene un entregable mayor y siete paquetes. Los seis de producto conservan el identificador y el monto del acta. El paquete 1.0 es la dirección del proyecto y se financia con la reserva de gestión. La contingencia no es un paquete: no produce un entregable.

\begin{table}[H]
\centering
\caption{Paquetes de la EDT. Los montos salen de la suma de actividades.}
\small
\begin{tabularx}{\textwidth}{@{}l X r@{}}
\toprule
\textbf{EDT} & \textbf{Paquete} & \textbf{USD} \\
\midrule
1.0 & Dirección del proyecto & 71\,000 \\
1.1 & Vista de posición del cliente & 480\,000 \\
1.2 & Pasarela de herramientas y consentimiento & 360\,000 \\
1.3 & Adaptadores de Banca Móvil y web & 280\,000 \\
1.4 & Descubrimiento con Yape, sin construirlo & 80\,000 \\
1.5 & Seguridad, conducta de mercado y datos & 140\,000 \\
1.6 & Habilitación de canales y del equipo & 80\,000 \\
\midrule
 & CAPEX base, paquetes 1.1 a 1.6 & 1\,420\,000 \\
 & Reserva de contingencia, 10\,\% del CAPEX base & 142\,000 \\
 & \textbf{Presupuesto autorizado} & \textbf{1\,633\,000} \\
\bottomrule
\end{tabularx}
\end{table}

La EDT cubre el trabajo autorizado. Queda fuera, y por eso no tiene código: desembolsos del agente, asesoría, rediseño de canales, reemplazo de Clara o de Salvador, escritura en los sistemas de producto y entrenar el modelo con conversaciones.

\subsection{Diccionario}

\begin{table}[H]
\centering
\caption{Diccionario de la EDT.}
\footnotesize
\begin{tabularx}{\textwidth}{@{}l X X@{}}
\toprule
\textbf{EDT} & \textbf{Entregable y aceptación} & \textbf{Cuándo y quién} \\
\midrule
1.0 & Plan de dirección firmado. Al cierre: traspaso, soporte, lecciones y la decisión sobre Yape. & H0, semana~4, y H6, semana~44. Director del proyecto. \\
1.1 & Vista de prueba para cuentas, tarjetas y créditos, con marca de vigencia. Si el dato no está vigente, la herramienta no afirma el saldo. Los sistemas de producto no se modifican. & H2, semana~16, mes~4. CIO. \\
1.2 & Seis habilidades, tres perfiles y revocación en la llamada siguiente. Simular no desembolsa. & Consulta en H3, Planificación en H4, Ofertas en H5. Banca Personas Digital. \\
1.3 & El agente se abre con la sesión que el cliente ya tiene, en la app y en la web. No hay una aplicación nueva. & H3, H4 y H5. Dueño de canales. \\
1.4 & Informe de enlace profundo o de una API que el dueño de Yape acepte. No se escribe en sus sistemas. & H6. Director del proyecto. \\
1.5 & Trazas, mínimo privilegio, vía humana y doce casos de comportamiento aprobados, sin efectos sobre saldos, pagos ni desembolsos. & Antes de H3 y antes de H5. CISO y Conducta de mercado. \\
1.6 & El personal de canales conoce el límite: informa y simula, no asesora ni coloca. & H3 y H5. Banca Personas Digital. \\
\bottomrule
\end{tabularx}
\end{table}

Los doce casos de prueba no forman un paquete aparte. El acta los pide como criterio de éxito y su costo está dentro de los \usd{140\,000} de seguridad y conducta: \usd{60\,000} para diseñarlos, \usd{20\,000} para el piloto y \usd{60\,000} para la salida a producción.

\section{Cronograma}

La unidad es la semana. El mes tiene cuatro semanas: el mes~1 cierra en la semana~4 y el mes~11 en la semana~44. Los hitos quedan en el mes previsto en el acta. H0 en el mes~1, H1 en el mes~2, H2 en el mes~4, H3 en el mes~6, H4 en el mes~8, H5 en el mes~10 y H6 en el mes~11.

La relación es fin-inicio, salvo el trabajo en paralelo. El inicio temprano es el fin de la predecesora más una semana. La holgura total es inicio tardío menos inicio temprano. La ruta crítica es la de holgura cero. El fin tardío del proyecto se ancla al fin temprano: si una actividad crítica se alarga, el cierre se mueve.

Las duraciones son juicio del equipo, por analogía con una integración bancaria de solo lectura. No hay histórico interno de un agente de este tipo, ni una estimación de tres puntos. Esas duraciones, y el costo de cada actividad, son datos de entrada del libro de trabajo.

\begin{table}[H]
\centering
\caption{Actividades, relaciones y holgura. En negrita, la ruta crítica.}
\footnotesize
\setlength{\tabcolsep}{3.5pt}
\begin{tabularx}{\textwidth}{@{}l l X l r r r@{}}
\toprule
\textbf{Id} & \textbf{EDT} & \textbf{Actividad} & \textbf{Predecesora} & \textbf{Duración} & \textbf{Fin} & \textbf{Holgura} \\
\midrule
\textbf{A01} & 1.0 & Inicio y plan de dirección & --- & 4 & 4 & 0 \\
\textbf{A02} & 1.0 & Arquitectura, amenazas y consentimiento & A01 & 4 & 8 & 0 \\
\textbf{A03} & 1.1 & Vista de posición en prueba & A02 & 8 & 16 & 0 \\
\textbf{A04} & 1.2 & Pasarela y consentimiento de Consulta & A02 & 8 & 16 & 0 \\
A05 & 1.5 & Diseño de controles y de los doce casos & A02 & 6 & 14 & 4 \\
A06 & 1.4 & Descubrimiento con Yape & A02 & 8 & 16 & 24 \\
\textbf{A07} & 1.3 & Adaptador de Consulta & A03, A04 & 6 & 22 & 0 \\
A08 & 1.6 & Habilitación del piloto & A05 & 4 & 18 & 4 \\
\textbf{A09} & 1.5 & Pruebas y piloto de Consulta & A07, A08, A05 & 2 & 24 & 0 \\
\textbf{A10} & 1.2 & Herramientas de Planificación & A09 & 4 & 28 & 0 \\
\textbf{A11} & 1.3 & Adaptador de Planificación & A10 & 4 & 32 & 0 \\
\textbf{A12} & 1.2 & Herramientas de Ofertas & A11 & 4 & 36 & 0 \\
\textbf{A13} & 1.3 & Adaptador de Ofertas & A11 & 4 & 36 & 0 \\
\textbf{A14} & 1.6 & Habilitación de la salida & A11 & 4 & 36 & 0 \\
\textbf{A15} & 1.5 & Pruebas de producción y tres firmas & A12, A13, A14 & 4 & 40 & 0 \\
\textbf{A16} & 1.0 & Cierre, traspaso y decisión sobre Yape & A15, A06 & 4 & 44 & 0 \\
\bottomrule
\end{tabularx}
\end{table}

\noindent La columna Fin es la semana de fin temprano. La dirección, \usd{16\,000}, corre de la semana~1 a la~44 y no entra en la red.

\subsection{Hitos}

\begin{table}[H]
\centering
\caption{Hitos de la línea base.}
\small
\begin{tabularx}{\textwidth}{@{}l c c X@{}}
\toprule
\textbf{Hito} & \textbf{Semana} & \textbf{Mes} & \textbf{Queda cerrado cuando} \\
\midrule
H0 & 4 & 1 & El plan de dirección está aprobado. \\
H1 & 8 & 2 & Hay arquitectura, análisis de amenazas y diseño del consentimiento. \\
H2 & 16 & 4 & La vista de cuentas, tarjetas y créditos responde en prueba. \\
H3 & 24 & 6 & El piloto de Consulta está en marcha con 5\,000 clientes y los doce casos pasaron. \\
H4 & 32 & 8 & Planificación está en la aplicación y en la web. \\
H5 & 40 & 10 & Ofertas con confirmación está en producción y hay vía humana. Conducta de mercado, CISO y Riesgo de crédito firmaron. \\
H6 & 44 & 11 & Hay traspaso, soporte y una decisión registrada sobre Yape. \\
\bottomrule
\end{tabularx}
\end{table}

\subsection{Ruta crítica}

La ruta crítica dura 44 semanas y pone los tres perfiles en producción:

A01, A02, A03 y A04 en paralelo, A07, A09, A10, A11, y en paralelo A12, A13 y A14, luego A15 y A16.

A03 y A04 terminan juntas en la semana~16. Si la vista de posición o la pasarela se atrasan, el piloto se atrasa. Es la dependencia que el acta ya identificaba a cargo del CIO. A12, A13 y A14 ocupan la misma ventana, de la semana~33 a la~36. Las tres tienen que terminar para que empiecen las pruebas de producción. Alargar cualquiera mueve H5 y H6.

Fuera de la ruta quedan tres actividades. El diseño de los doce casos (A05) y la habilitación del piloto (A08) pueden demorarse cuatro semanas. El descubrimiento con Yape (A06) puede demorarse 24 semanas: su única sucesora es el cierre, que espera a la semana~41. Por eso Yape es Should y no Must.

En la red, la flecha es fin-inicio y el recorrido va de izquierda a derecha: A01 abre el proyecto y A16 lo cierra. A03 y A04 son las dos críticas, igual que A12, A13 y A14: cualquiera de esas ramas, sumada al resto, da 44 semanas.

\begin{landscape}
\thispagestyle{empty}
\noindent\textcolor{bcp}{\small Agente financiero personal}\hfill\textcolor{bcp}{\small Línea base de alcance, cronograma y costo}\\[-0.35em]
{\color{cobre}\rule{\linewidth}{0.4pt}}\par
\vspace{0.6em}
\begin{figure}[H]
\centering
\caption{Red del cronograma. En cobre, la ruta crítica. En azul, las actividades con holgura.}
\label{fig:red}
\begin{tikzpicture}[
  x=1cm,y=1cm,
  act/.style={rectangle, rounded corners=1.5pt, text width=1.78cm, minimum height=1.42cm, align=center, inner sep=2.5pt, font=\scriptsize, text=white},
  crit/.style={act, fill=cobre, draw=cobre!55!black, line width=0.6pt},
  slack/.style={act, fill=bcp, draw=bcp!75!black, line width=0.6pt},
  term/.style={act, rounded corners=0.72cm, fill=cobre, draw=cobre!40!black, line width=1.2pt},
  arr/.style={-{Latex[length=2.1mm,width=1.7mm]}, draw=cobre, line width=0.85pt},
  arrs/.style={-{Latex[length=2.1mm,width=1.7mm]}, draw=bcp, line width=0.7pt},
  plain/.style={draw=cobre, line width=0.85pt},
  plains/.style={draw=bcp, line width=0.7pt}
]
  \node[term] (A01) at (1.12,0) {{\footnotesize\bfseries INICIO}\\[1pt] A01 · plan\\4 sem};
  \node[crit] (A02) at (3.52,0) {{\bfseries A02}\\Arquitectura\\4 sem};
  \node[crit] (A03) at (5.92,1.82) {{\bfseries A03}\\Vista\\8 sem};
  \node[crit] (A04) at (5.92,-1.82) {{\bfseries A04}\\Pasarela\\8 sem};
  \node[slack] (A05) at (5.92,-3.72) {{\bfseries A05}\\Doce casos\\holgura 4};
  \node[slack] (A06) at (5.92,-5.55) {{\bfseries A06}\\Yape\\holgura 24};
  \node[crit] (A07) at (8.32,0) {{\bfseries A07}\\Consulta\\6 sem};
  \node[slack] (A08) at (8.32,-3.72) {{\bfseries A08}\\Habilitación\\holgura 4};
  \node[crit] (A09) at (10.72,0) {{\bfseries A09}\\Piloto\\2 sem};
  \node[crit] (A10) at (13.12,0) {{\bfseries A10}\\Planificación\\4 sem};
  \node[crit] (A11) at (15.52,0) {{\bfseries A11}\\Adapta plan\\4 sem};
  \node[crit] (A12) at (17.92,1.82) {{\bfseries A12}\\Ofertas\\4 sem};
  \node[crit] (A13) at (17.92,0) {{\bfseries A13}\\Adapta ofertas\\4 sem};
  \node[crit] (A14) at (17.92,-1.82) {{\bfseries A14}\\Salida\\4 sem};
  \node[crit] (A15) at (20.32,0) {{\bfseries A15}\\Producción\\4 sem};
  \node[term, line width=1.35pt] (A16) at (22.72,0) {{\footnotesize\bfseries FIN}\\[1pt] A16 · cierre\\semana 44};

  \node[font=\scriptsize\bfseries, text=cobre, anchor=south] at (1.12,0.82) {Semana 1};
  \node[font=\scriptsize\bfseries, text=cobre, anchor=south] at (22.72,0.82) {Semana 44};

  \draw[arr] (A01) -- (A02);
  \draw[plain] (A02.east) -- (4.78,0);
  \draw[arr] (4.78,0) -- (4.78,1.82) -- (A03.west);
  \draw[arr] (4.78,0) -- (4.78,-1.82) -- (A04.west);
  \draw[arr] (A03.east) -- (7.18,1.82) -- (7.18,0.26) -- ([yshift=2.4mm]A07.west);
  \draw[arr] (A04.east) -- (7.18,-1.82) -- (7.18,-0.26) -- ([yshift=-2.4mm]A07.west);
  \draw[arr] (A07) -- (A09);
  \draw[arr] (A09) -- (A10);
  \draw[arr] (A10) -- (A11);
  \draw[plain] (A11.east) -- (16.78,0);
  \draw[arr] (16.78,0) -- (16.78,1.82) -- (A12.west);
  \draw[arr] (16.78,0) -- (A13.west);
  \draw[arr] (16.78,0) -- (16.78,-1.82) -- (A14.west);
  \draw[arr] (A12.east) -- (19.18,1.82) -- (19.18,0.26) -- ([yshift=2.4mm]A15.west);
  \draw[arr] (A13) -- (A15);
  \draw[arr] (A14.east) -- (19.18,-1.82) -- (19.18,-0.26) -- ([yshift=-2.4mm]A15.west);
  \draw[arr] (A15) -- (A16);

  \draw[plains] (A02.south) -- (3.52,-3.72);
  \draw[arrs] (3.52,-3.72) -- (A05.west);
  \draw[plains] (3.52,-3.72) -- (3.52,-5.55);
  \draw[arrs] (3.52,-5.55) -- (A06.west);
  \draw[arrs] (A05) -- (A08);
  \draw[arrs] (A08.east) -- (10.72,-3.72) -- (A09.south);
  \draw[arrs] (A06.east) -- (22.72,-5.55) -- (A16.south);

  \node[crit, minimum width=0.38cm, minimum height=0.24cm, text width=0cm, inner sep=0] (lg1) at (0.35,-6.85) {};
  \node[anchor=west, font=\scriptsize, text=tinta] at (0.62,-6.85) {Ruta crítica, holgura cero};
  \node[slack, minimum width=0.38cm, minimum height=0.24cm, text width=0cm, inner sep=0] (lg2) at (5.55,-6.85) {};
  \node[anchor=west, font=\scriptsize, text=tinta] at (5.82,-6.85) {Con holgura};
  \draw[arr] (9.15,-6.85) -- (9.85,-6.85);
  \node[anchor=west, font=\scriptsize, text=tinta] at (9.98,-6.85) {Fin-inicio};
\end{tikzpicture}
\end{figure}
\vfill
\noindent\textcolor{tinta}{\footnotesize Banca Personas · BCP}\hfill\textcolor{tinta}{\footnotesize\thepage}
\end{landscape}

\begin{figure}[H]
\centering
\caption{Cronograma de 44 semanas. La barra cobre es ruta crítica. La barra azul tiene holgura, y la franja clara mide esa holgura.}
\label{fig:gantt}
\begin{tikzpicture}[x=0.242cm,y=0.46cm]
  \foreach \x in {0,4,8,12,16,20,24,28,32,36,40,44}
    \draw[black!15] (\x,0.35) -- (\x,16.7);
  \foreach \m/\lab in {2/1,6/2,10/3,14/4,18/5,22/6,26/7,30/8,34/9,38/10,42/11}
    \node[font=\scriptsize,text=bcp] at (\m,17.15) {\lab};
  \node[anchor=east,font=\scriptsize,text=bcp] at (-0.45,17.15) {Mes};
  % holgura
  \fill[suave] (14,11.72) rectangle (18,12.28);
  \fill[suave] (16,10.72) rectangle (40,11.28);
  \fill[suave] (18,8.72) rectangle (22,9.28);
  % ruta crítica
  \fill[cobre] (0,15.72) rectangle (4,16.28);
  \fill[cobre] (4,14.72) rectangle (8,15.28);
  \fill[cobre] (8,13.72) rectangle (16,14.28);
  \fill[cobre] (8,12.72) rectangle (16,13.28);
  \fill[cobre] (16,9.72) rectangle (22,10.28);
  \fill[cobre] (22,7.72) rectangle (24,8.28);
  \fill[cobre] (24,6.72) rectangle (28,7.28);
  \fill[cobre] (28,5.72) rectangle (32,6.28);
  \fill[cobre] (32,4.72) rectangle (36,5.28);
  \fill[cobre] (32,3.72) rectangle (36,4.28);
  \fill[cobre] (32,2.72) rectangle (36,3.28);
  \fill[cobre] (36,1.72) rectangle (40,2.28);
  \fill[cobre] (40,0.72) rectangle (44,1.28);
  % con holgura
  \fill[bcp] (8,11.72) rectangle (14,12.28);
  \fill[bcp] (8,10.72) rectangle (16,11.28);
  \fill[bcp] (14,8.72) rectangle (18,9.28);
  \node[anchor=east,font=\scriptsize] at (-0.4,16) {A01 Inicio};
  \node[anchor=east,font=\scriptsize] at (-0.4,15) {A02 Arquitectura};
  \node[anchor=east,font=\scriptsize] at (-0.4,14) {A03 Vista de posición};
  \node[anchor=east,font=\scriptsize] at (-0.4,13) {A04 Pasarela Consulta};
  \node[anchor=east,font=\scriptsize] at (-0.4,12) {A05 Diseño de pruebas};
  \node[anchor=east,font=\scriptsize] at (-0.4,11) {A06 Descubrimiento Yape};
  \node[anchor=east,font=\scriptsize] at (-0.4,10) {A07 Adaptador Consulta};
  \node[anchor=east,font=\scriptsize] at (-0.4,9) {A08 Habilitación piloto};
  \node[anchor=east,font=\scriptsize] at (-0.4,8) {A09 Piloto Consulta};
  \node[anchor=east,font=\scriptsize] at (-0.4,7) {A10 Planificación};
  \node[anchor=east,font=\scriptsize] at (-0.4,6) {A11 Adaptador Planificación};
  \node[anchor=east,font=\scriptsize] at (-0.4,5) {A12 Ofertas};
  \node[anchor=east,font=\scriptsize] at (-0.4,4) {A13 Adaptador Ofertas};
  \node[anchor=east,font=\scriptsize] at (-0.4,3) {A14 Habilitación salida};
  \node[anchor=east,font=\scriptsize] at (-0.4,2) {A15 Pruebas de producción};
  \node[anchor=east,font=\scriptsize] at (-0.4,1) {A16 Cierre};
  \foreach \x/\h in {4/H0,8/H1,16/H2,24/H3,32/H4,40/H5,44/H6}{
    \node[diamond,fill=cobre,inner sep=1.15pt] at (\x,0.22) {};
    \node[anchor=north,font=\scriptsize] at (\x,-0.02) {\h};
  }
\end{tikzpicture}
\end{figure}

\section{Presupuesto, recursos y supuestos}

\subsection{De la actividad al presupuesto}

El costo de un paquete es la suma de sus actividades. El acta fijó el monto del paquete; aquí solo se reparte entre actividades, según el peso del trabajo. No hay tarifa por hora. La vista de posición queda en una sola actividad, porque es un solo entregable de prueba. La pasarela y los adaptadores se reparten entre los tres perfiles, con más peso en Consulta, que abre el camino.

\begin{table}[H]
\centering
\caption{Costo por actividad. Suma de 1.1 a 1.6: \usd{1\,420\,000}. Suma de 1.0: \usd{71\,000}.}
\footnotesize
\setlength{\tabcolsep}{4pt}
\begin{tabularx}{\textwidth}{@{}l l X r@{}}
\toprule
\textbf{Id} & \textbf{EDT} & \textbf{Actividad} & \textbf{USD} \\
\midrule
A01 & 1.0 & Inicio y plan de dirección & 15\,000 \\
A02 & 1.0 & Arquitectura, amenazas y consentimiento & 20\,000 \\
A16 & 1.0 & Cierre, traspaso y decisión sobre Yape & 20\,000 \\
--- & 1.0 & Dirección durante los 11 meses & 16\,000 \\
A03 & 1.1 & Vista de posición en prueba & 480\,000 \\
A04 & 1.2 & Pasarela de Consulta & 200\,000 \\
A10 & 1.2 & Herramientas de Planificación & 80\,000 \\
A12 & 1.2 & Herramientas de Ofertas & 80\,000 \\
A07 & 1.3 & Adaptador de Consulta & 110\,000 \\
A11 & 1.3 & Adaptador de Planificación & 90\,000 \\
A13 & 1.3 & Adaptador de Ofertas & 80\,000 \\
A06 & 1.4 & Descubrimiento con Yape & 80\,000 \\
A05 & 1.5 & Diseño de controles y de los doce casos & 60\,000 \\
A09 & 1.5 & Pruebas y piloto de Consulta & 20\,000 \\
A15 & 1.5 & Pruebas de producción y tres firmas & 60\,000 \\
A08 & 1.6 & Habilitación del piloto & 30\,000 \\
A14 & 1.6 & Habilitación de la salida & 50\,000 \\
\bottomrule
\end{tabularx}
\end{table}

La contingencia es \usd{142\,000}, el 10\,\% del CAPEX base, y no se asigna a una actividad. Cubre dos riesgos del acta, ambos con probabilidad~3, impacto~5 y severidad~15. El porcentaje es el del acta. No es un valor monetario esperado, y los \usd{142\,000} son una sola reserva para los dos riesgos.

\begin{table}[H]
\centering
\caption{Riesgos que puede cubrir la contingencia. La reserva sigue sin asignar.}
\small
\begin{tabularx}{\textwidth}{@{}l X l X@{}}
\toprule
\textbf{Id} & \textbf{Riesgo} & \textbf{Paquete} & \textbf{Para qué se usaría} \\
\midrule
R01 & El agente responde con una cifra de una vista vencida. & 1.1 y 1.5 & Rehacer la marca de vigencia o repetir las pruebas. \\
R07 & El CIO atrasa la vista de posición y H2 pasa de la semana~16. & 1.1 & Sostener la vista hasta que responda en prueba. \\
\bottomrule
\end{tabularx}
\end{table}

El caso de negocio supone que esa reserva se gasta. En esta línea base permanece sin asignar hasta que el riesgo ocurra. El director del proyecto puede usarla hasta el límite del acta, sin una aprobación nueva.

El presupuesto autorizado es CAPEX base más dirección más contingencia: \(1\,420\,000 + 71\,000 + 142\,000 = 1\,633\,000\).

\begin{figure}[H]
\centering
\caption{Presupuesto autorizado, en miles de USD. La contingencia no es un paquete: queda aparte hasta que ocurra el riesgo.}
\label{fig:presupuesto}
\begin{tikzpicture}[x=0.021cm,y=0.72cm]
  \foreach \v in {0,100,200,300,400,480}
    \draw[black!12] (\v,0.2) -- (\v,8.5);
  \foreach \v/\lab in {0/0,100/100,200/200,300/300,400/400,480/480}
    \node[font=\scriptsize,text=bcp,below] at (\v,0.15) {\lab};
  \fill[bcp] (0,7.25) rectangle (480,7.85);
  \fill[bcp] (0,6.25) rectangle (360,6.85);
  \fill[bcp] (0,5.25) rectangle (280,5.85);
  \fill[cobre] (0,4.25) rectangle (142,4.85);
  \fill[bcp] (0,3.25) rectangle (140,3.85);
  \fill[bcp] (0,2.25) rectangle (80,2.85);
  \fill[bcp] (0,1.25) rectangle (80,1.85);
  \fill[bcp] (0,0.25) rectangle (71,0.85);
  \node[anchor=east,font=\scriptsize] at (-8,7.55) {1.1 Vista};
  \node[anchor=east,font=\scriptsize] at (-8,6.55) {1.2 Pasarela};
  \node[anchor=east,font=\scriptsize] at (-8,5.55) {1.3 Adaptadores};
  \node[anchor=east,font=\scriptsize] at (-8,4.55) {Contingencia};
  \node[anchor=east,font=\scriptsize] at (-8,3.55) {1.5 Seguridad};
  \node[anchor=east,font=\scriptsize] at (-8,2.55) {1.4 Yape};
  \node[anchor=east,font=\scriptsize] at (-8,1.55) {1.6 Habilitación};
  \node[anchor=east,font=\scriptsize] at (-8,0.55) {1.0 Dirección};
  \node[anchor=west,font=\scriptsize] at (488,7.55) {480};
  \node[anchor=west,font=\scriptsize] at (368,6.55) {360};
  \node[anchor=west,font=\scriptsize] at (288,5.55) {280};
  \node[anchor=west,font=\scriptsize] at (150,4.55) {142};
  \node[anchor=west,font=\scriptsize] at (148,3.55) {140};
  \node[anchor=west,font=\scriptsize] at (88,2.55) {80};
  \node[anchor=west,font=\scriptsize] at (88,1.55) {80};
  \node[anchor=west,font=\scriptsize] at (79,0.55) {71};
  \node[font=\scriptsize,text=bcp] at (240,-0.85) {Miles de USD};
\end{tikzpicture}
\end{figure}

\subsection{Gasto por mes}

El gasto de cada actividad se reparte entre las semanas que van de su inicio temprano a su fin temprano. Si el costo no es divisible entre esas semanas, el residuo de dólares queda en las últimas. La contingencia no entra en la tabla: no tiene mes hasta que un riesgo la active. El descubrimiento con Yape aparece en los meses~3 y~4 porque esa es su posición temprana; las 24 semanas de holgura permiten correr esos \usd{80\,000} hasta los meses~9 y~10 sin mover el cierre.

\begin{table}[H]
\centering
\caption{Gasto de las actividades en los meses 1 a 6, en USD.}
\footnotesize
\setlength{\tabcolsep}{4pt}
\begin{tabularx}{\textwidth}{@{}l *{6}{>{\raggedleft\arraybackslash}X}@{}}
\toprule
\textbf{EDT} & \textbf{1} & \textbf{2} & \textbf{3} & \textbf{4} & \textbf{5} & \textbf{6} \\
\midrule
1.0 & 16\,452 & 21\,452 & 1\,452 & 1\,452 & 1\,456 & 1\,456 \\
1.1 & 0 & 0 & 240\,000 & 240\,000 & 0 & 0 \\
1.2 & 0 & 0 & 100\,000 & 100\,000 & 0 & 0 \\
1.3 & 0 & 0 & 0 & 0 & 73\,332 & 36\,668 \\
1.4 & 0 & 0 & 40\,000 & 40\,000 & 0 & 0 \\
1.5 & 0 & 0 & 40\,000 & 20\,000 & 0 & 20\,000 \\
1.6 & 0 & 0 & 0 & 15\,000 & 15\,000 & 0 \\
\midrule
Mes & 16\,452 & 21\,452 & 421\,452 & 416\,452 & 89\,788 & 58\,124 \\
\bottomrule
\end{tabularx}
\end{table}

\begin{table}[H]
\centering
\caption{Gasto de las actividades en los meses 7 a 11, en USD. La última columna es el total del paquete.}
\footnotesize
\setlength{\tabcolsep}{4pt}
\begin{tabularx}{\textwidth}{@{}l *{6}{>{\raggedleft\arraybackslash}X}@{}}
\toprule
\textbf{EDT} & \textbf{7} & \textbf{8} & \textbf{9} & \textbf{10} & \textbf{11} & \textbf{Total} \\
\midrule
1.0 & 1\,456 & 1\,456 & 1\,456 & 1\,456 & 21\,456 & 71\,000 \\
1.1 & 0 & 0 & 0 & 0 & 0 & 480\,000 \\
1.2 & 80\,000 & 0 & 80\,000 & 0 & 0 & 360\,000 \\
1.3 & 0 & 90\,000 & 80\,000 & 0 & 0 & 280\,000 \\
1.4 & 0 & 0 & 0 & 0 & 0 & 80\,000 \\
1.5 & 0 & 0 & 0 & 60\,000 & 0 & 140\,000 \\
1.6 & 0 & 0 & 50\,000 & 0 & 0 & 80\,000 \\
\midrule
Mes & 81\,456 & 91\,456 & 211\,456 & 61\,456 & 21\,456 & 1\,491\,000 \\
\bottomrule
\end{tabularx}
\end{table}

Al cierre del mes~4 van \usd{875\,808}, el 58,7\,\% del gasto de las actividades, porque en ese mes terminan la vista de posición y la pasarela de Consulta. Al mes~11 el acumulado es \usd{1\,491\,000}: CAPEX base más dirección. Los \usd{142\,000} de contingencia no tienen mes, y el presupuesto autorizado sigue en \usd{1\,633\,000}. Cuando el costo de una actividad no es divisible entre sus semanas, el residuo queda en las últimas. De ahí salen los \usd{73\,332} y \usd{36\,668} del adaptador de Consulta, y la diferencia de unos dólares en la dirección mensual.

\subsection{Recursos}

\begin{table}[H]
\centering
\caption{Responsable de cada paquete. El costo ya está en el paquete y no se recalcula por persona.}
\small
\begin{tabularx}{\textwidth}{@{}l X X@{}}
\toprule
\textbf{EDT} & \textbf{Quién hace el trabajo} & \textbf{Límite del trabajo} \\
\midrule
1.0 & Un director de proyecto durante los 11 meses. & No promete la conversión del 2\,\% ni compromete a Yape. \\
1.1 & Equipo del CIO y dueños de cuentas, tarjetas y créditos. & No modifica esos sistemas. \\
1.2 & Equipo de herramientas y consentimiento. & No abre el catálogo ni ejecuta un pago. \\
1.3 & Equipo de canales, sobre la sesión existente. & No rediseña la aplicación ni la banca por internet. \\
1.4 & Un analista de descubrimiento. & No escribe en Yape. \\
1.5 & Conducta de mercado, CISO y Riesgo de crédito. & No firman un despliegue si un caso falla. \\
1.6 & Equipo de canales, en formación. & No convierte al personal en colocadores del agente. \\
\bottomrule
\end{tabularx}
\end{table}

\subsection{Supuestos de esta línea base}

\begin{enumerate}
  \item El mes tiene cuatro semanas y el proyecto dura 44 semanas, del mes~1 al mes~11 del acta.
  \item Las relaciones son fin-inicio. Donde hay paralelo, las ramas tienen que terminar antes de la actividad que las reúne.
  \item La duración es juicio del equipo, por analogía con una integración bancaria de solo lectura. No hay histórico interno ni estimación de tres puntos. En el libro de trabajo es un dato de entrada.
  \item El costo del paquete es el del acta. El reparto entre actividades es del equipo y, en el libro de trabajo, es un dato de entrada.
  \item La contingencia es el 10\,\% del CAPEX base. No se compone con la reserva de gestión y no sale de un valor monetario esperado: cubre R01 y R07, y sigue sin asignar.
  \item El gasto mensual usa el inicio temprano. El residuo de dólares, cuando el costo no es divisible entre las semanas, queda en las últimas.
  \item El OPEX de cinco años, \usd{2\,217\,741}, empieza después del cierre. No entra en los \usd{1\,633\,000}.
  \item No se resta el sueldo de los asesores que hoy llevan cartera. El agente no toma ese trabajo: informa y simula, y la norma exige una vía humana.
\end{enumerate}

\section{Trazabilidad}

Cada paquete queda ligado a un hito y a un monto.

\begin{table}[H]
\centering
\caption{Del paquete al hito y al costo.}
\small
\begin{tabularx}{\textwidth}{@{}l X c r@{}}
\toprule
\textbf{EDT} & \textbf{Se acepta en} & \textbf{Semana} & \textbf{USD} \\
\midrule
1.0 & H0 el plan, H6 el cierre & 4 y 44 & 71\,000 \\
1.1 & H2, vista en prueba & 16 & 480\,000 \\
1.2 & H3, H4 y H5, un perfil por hito & 24, 32 y 40 & 360\,000 \\
1.3 & H3, H4 y H5, el adaptador de cada perfil & 24, 32 y 40 & 280\,000 \\
1.4 & H6, decisión sobre Yape & 44 & 80\,000 \\
1.5 & H3 el piloto y H5 la producción & 24 y 40 & 140\,000 \\
1.6 & H3 el piloto y H5 la salida & 24 y 40 & 80\,000 \\
\bottomrule
\end{tabularx}
\end{table}

Si cambia la duración de A03, la vista de posición, se mueven H2, el piloto y los hitos posteriores, y cambia la holgura de Yape. Si cambia el costo de una actividad, el total del paquete deja de coincidir con el monto del acta, y el presupuesto autorizado tampoco.

\end{document}
